\documentclass[main.tex]{subfiles}


\begin{document}

%from pagenumber to up
% \phantomsection 
% \hypertarget{toc}{}

 

 

% номер секции вручную
\setcounter{section}{7
}
\addtocounter{section}{-1} 

\section{Линии электрического поля}


Наглядно описывать электрическое поле в пространстве принято с помощью \textbf{линий электрического поля}~--- линий, идущих вдоль векторов напряженности поля в различных точках пространства\footnote{Эти линии называют также \emph{линиями напряженности} или \emph{силовыми линиями}. Всякая линия поля изображается так, что в любой ее точке вектор напряженности направлен по касательной к линии поля.}.

На рис.\,\ref{fig:p99} показаны картины линий поля \emph{точечного заряда} для двух случаев: положительный и отрицательный заряд.

    \begin{figure}[H]
    \centering

    \begin{tikzpicture}[font=\small,            line cap=round,                             >={Stealth[inset=0pt,round,length=3mm, angle'=20]},vec/.style={>={Stealth[inset=0pt,round,length=3.5mm, angle'=20]}}]


%line
\foreach \i in{0,...,11}
\draw[Green,decoration={markings,mark=at position {1/2+3/20} with {\arrow{>}}},postaction={decorate}] (0,0) --++ ({30*\i}:2);

\shade[ball color=red] (0,0) circle (.15);

%%%%%%%%%%%%%%%%%

\begin{scope}[xshift=6cm]

%line
\foreach \i in{0,...,11}
\draw[Green,decoration={markings,mark=at position {1/2} with {\arrowreversed{>}}},postaction={decorate}] (0,0) --++ ({30*\i}:2);

\shade[ball color=NavyBlue] (0,0) circle (.15);

    
\end{scope}



    \end{tikzpicture}
\caption{Линии электрического поля точечного заряда}
\label{fig:p99}
\end{figure}

Линии поля положительного заряда \emph{выходят из него} (рис.\,\ref{fig:p99}, слева); наоборот, линии поля отрицательного заряда \emph{входят в него} (рис.\,\ref{fig:p99}, справа).  

\emph{Линии поля указывают направления векторов напряженности в пространстве. Чем больше густота линий, тем больше напряженность в данной области пространства.}

На рис.\,\ref{fig:p100} показаны картины линий поля вблизи равномерно \emph{заряженной пластины} в двух случаях: положительно и отрицательно заряженная пластина.

    \begin{figure}[H]
    \centering

    \begin{tikzpicture}[font=\small,            line cap=round,                             >={Stealth[inset=0pt,round,length=3mm, angle'=20]},vec/.style={>={Stealth[inset=0pt,round,length=3.5mm, angle'=20]}}]


%line
\foreach \i in{0,...,5}
\draw[Green,decoration={markings,mark=at position {1/2+1.5/20} with {\arrow{>}}},postaction={decorate}] (0,{.25+\i*.5}) --++ ({180}:2);

\foreach \i in{0,...,5}
\draw[Green,decoration={markings,mark=at position {1/2+1.5/20} with {\arrow{>}}},postaction={decorate}] (0,{.25+\i*.5}) --++ ({0}:2);


%plane
\draw[thick,red] (0,0) --++ (0,3);

%%%%%%%%%%%%%%%%%%%%%%

\begin{scope}[xshift=6cm]

%line
\foreach \i in{0,...,5}
\draw[Green,decoration={markings,mark=at position {1/2+1.5/20} with {\arrow{<}}},postaction={decorate}] (0,{.25+\i*.5}) --++ ({180}:2);

\foreach \i in{0,...,5}
\draw[Green,decoration={markings,mark=at position {1/2+1.5/20} with {\arrow{<}}},postaction={decorate}] (0,{.25+\i*.5}) --++ ({0}:2);


%plane
\draw[thick,NavyBlue] (0,0) --++ (0,3);
    
\end{scope}


    \end{tikzpicture}
\caption{Линии электрического поля заряженной пластины}
\label{fig:p100}
\end{figure}

Возле каждой из пластин на рис.\,\ref{fig:p100} поле является \emph{однородным}. Из этого рисунка видно, что \emph{линии однородного поля являются параллельными прямыми одинаковой густоты.}  

\emph{Линии поля не пересекаются, они всегда начинаются на положительных зарядах и заканчиваются на отрицательных.}

\medskip

\noindent\textbf{Задача.} Как меняется напряженность поля с удалением от точечного заряда?

\smallskip

\noindent\emph{Решение}. С увеличением расстояния от заряда, создающего поле, густота линий поля уменьшается (см.~рис.\,\ref{fig:p99}); значит, с удалением от заряда напряженность уменьшается.




 
\end{document}